% !PRINTED 39
\weberpage{39}
At the same time, as formula (17) shows, if $\Theta_1$ is an arbitrary element in $S$, then $\Theta_1^{a_\nu}$ is contained in $S_{\nu-1}$, and the exponents $\lambda_1, \lambda_2 \ldots \lambda_{\nu-1}$ are divisible by $a_\nu$. Hence requirement 2) is also satisfied.

With this, then, our Theorem I. is proved. At the same time we see from this derivation \gesperrt{that for an arbitrary Abelian group $S$ the elements of the basis can always be taken in such a way that their degrees are prime powers.}

From the representability by a basis it also follows that every Abelian group is metacyclic; for if the elements of such a group $S$ are represented in the form:
\[
A = A_1^{\alpha_1} A_2^{\alpha_2} \ldots A_\nu^{\alpha_\nu},
\]
and if $p$ is a prime which divides the degree $a_1$ of $A_1$, then the elements
\[
A' = A_1^{\alpha_1 p} A_2^{\alpha_2} \ldots A_\nu^{\alpha_\nu},
\]
when $\alpha_1$ runs through a complete system of residues to the modulus $a_1 : p$, form a divisor $S'$ of $S$ of index $p$, which is representable by the basis $A_1^p, A_2 \ldots A_\nu$. From $S'$ one can again find in the same way a divisor of prime index, and so forth. Hence $S$ is metacyclic (\S.\,8).

\begin{center}
\S.\,10.

\gesperrt{The Invariants of the Abelian Groups.}
\end{center}

The proof which we communicated in the previous section for the possibility of representing an Abelian group by a basis contains at the same time a way of finding such a basis, and indeed one in which the degrees of the basis elements are prime powers. Nevertheless it can happen that one and the same group can be represented by bases in different ways.

If we consider, for example, two elements $A, B$ whose degrees $a, b$ are relatively prime, then by these, as elements of a two-term basis, a group
\weq{1}{A^\alpha B^\beta \qquad \begin{aligned}
\alpha &= 0, 1, 2 \ldots a-1 \\
\beta &= 0, 1, 2 \ldots b-1
\end{aligned}}
is represented. But the same group can also be represented by the one-term basis $AB$ in the form
\weq{2}{(AB)^s \qquad s = 0, 1, 2 \ldots ab-1.}
